\chapter{Visualization and interactive tools}
\label{ch:visualization}

ATTPCROOTv2 has two kinds of visual tools. The first is the compiled event viewer in
\code{AtEventDisplay/}, built on ROOT's TEve, which browses the events of a FairRoot output file
and can re-run FairRoot tasks on the displayed event. The second is a set of interactive ROOT
macros and one compiled class (\code{AtUKFDisplay}) that serve a narrower purpose: refitting a
single track with live parameters, stepping through fitted events, and drawing particle-identification
(PID) gates. This chapter describes both, starting from the viewer.

\paragraph{Display requirement.} Every tool in this chapter that opens a window (the TEve
viewer, the fitter GUIs, the gate drawers) needs an X display. The TEve viewer additionally needs
working OpenGL; several macro headers in the a1975 workspace state that the Eve/OpenGL viewer does
not work under WSL (for example \url{macro/Unpack_HDF5/a1975/UKF/pipeline/event_png_UKF.C}). For
headless use, run ROOT with \code{root -l -b -q} and save canvases to files with
\code{TCanvas::SaveAs}; Section~\ref{sec:vis-batch} lists the macros written for that purpose.
Running an interactive tool with \code{-b} does not work: the gate drawers state explicitly that
they must be started with \code{root -l}, never \code{-b}.

\section{The event viewer: \code{AtViewerManager}}
\label{sec:vis-viewer}

\code{AtViewerManager} (\url{AtEventDisplay/AtViewerManager.h}) is the controller of the current
viewer. It is a singleton: the constructor stores \code{this} in a static pointer, and a second
construction aborts with \code{LOG(fatal)} (``Attempting to create a second instance of
AtViewerManager''). \code{AtViewerManager::Instance()} returns the pointer after construction.
The class is built in the \code{AtEventDisplay} library (\code{libAtEventDisplay.so}).

\subsection{What the constructor and \code{Init()} do}

The constructor takes a \code{std::shared\_ptr<AtMap>}, the pad map of the experiment, and builds
the window:
\begin{enumerate}
  \item it calls \code{TEveManager::Create(true, "IV")}, which creates a mapped TEve manager
        without the file browser;
  \item it embeds a vertical \code{AtEventSidebar} in the left panel of the TEve browser, titled
        ``Control'', holding \code{AtSidebarRunInfo} (input file, run ID, number of entries) and
        \code{AtSidebarPadControl} (the selected pad);
  \item it embeds a horizontal \code{AtEventSidebar} (the ``basebar'') in the bottom panel,
        holding \code{AtSidebarEventControl} (event number, Prev/Next, ``Rerun Event'') and
        \code{AtSidebarBranchControl} (three drop-down menus for the branch types);
  \item it attaches itself as an observer of the current tree entry.
\end{enumerate}

\code{Init()} sets the palette with \code{gStyle->SetPalette(55)}, calls the
FairRoot initialisation \code{FairRunAna::Instance()->Init()}, initialises every tab, builds the list of branches present
in the input, fills the sidebar and basebar widgets, and loads entry 0. Because
\code{AddTask()} and \code{Init()} both go through \code{FairRunAna::Instance()}, the macro must
create and configure a \code{FairRunAna} (source, sink, geometry file, runtime database)
\emph{before} constructing the viewer. Tabs must also be created after the viewer: the
constructors of \code{AtTabMain}, \code{AtTabPad} and \code{AtTabMacro} throw the string
``AtViewerManager must be initialized before creating tabs!'' if \code{Instance()} is null.

\subsection{Public interface}

\begin{center}
\begin{tabular}{p{0.38\textwidth}p{0.55\textwidth}}
\toprule
Method & Effect \\
\midrule
Constructor, argument \code{shared\_ptr<AtMap>} & Creates the TEve window, sidebar and basebar. \\
\code{AddTab(unique\_ptr<AtTabBase>)} & Takes ownership of a tab; tabs appear in the order added. \\
\code{AddTask(FairTask*)} & Forwards to \code{FairRunAna::Instance()->AddTask}. The task runs
  every time an entry is loaded. \\
\code{Init()} & Initialises the run, the tabs and the GUI; goes to entry 0. \\
\code{GotoEvent(Int\_t)}, \code{NextEvent()}, \code{PrevEvent()} & Navigation. \\
\code{SetCheckBranch(AtBranch\&)} & Makes Next/Prev skip entries whose \code{AtRawEvent} in that
  branch is not \code{IsGood()}. \\
\code{GetSidebar()} & The left sidebar, to which add-ons are added. \\
\code{GetMap()}, \code{GetBranchNames()} & The pad map and the type-to-branch table. \\
\code{GetRawEventBranch()}, \code{GetEventBranch()}, \code{GetPatternEventBranch()},
\code{GetCurrentEntry()}, \code{GetPadNum()} & The shared subjects described next. \\
\bottomrule
\end{tabular}
\end{center}

\section{The subject/observer design}
\label{sec:vis-observer}

The viewer does not pass data between widgets directly. Every piece of shared state is a
\emph{subject} (\code{DataHandling::AtSubject}, \url{AtData/AtDataSubject.h}), and every widget
or tab that depends on it is an \emph{observer} (\code{DataHandling::AtObserver},
\url{AtData/AtDataObserver.h}) that registers with \code{Attach()} and is called back through
\code{Update(AtSubject*)} when the subject calls \code{Notify()}.

\code{AtViewerManager} owns five subjects:
\begin{center}
\begin{tabular}{p{0.26\textwidth}p{0.22\textwidth}p{0.42\textwidth}}
\toprule
Member & Type & Meaning \\
\midrule
\code{fEntry} & \code{AtTreeEntry} & Current entry of the FairRoot tree. \\
\code{fRawEventBranch} & \code{AtBranch} & Branch shown as the \code{AtRawEvent}. \\
\code{fEventBranch} & \code{AtBranch} & Branch shown as the \code{AtEvent} (hits). \\
\code{fPatternEventBranch} & \code{AtBranch} & Branch shown as the \code{AtPatternEvent}. \\
\code{fPadNum} & \code{AtPadNum} = \code{AtSimpleType<int>} & Currently selected pad, initial
  value $-1$. \\
\bottomrule
\end{tabular}
\end{center}

The subjects carry behaviour, not only data (\url{AtEventDisplay/AtViewerManagerSubject.cxx}):
\begin{itemize}
  \item \code{AtTreeEntry::Set(entry)} calls \code{FairRunAna::Instance()->Run(entry)} and only
        then notifies. Changing the entry therefore re-executes every FairTask added with
        \code{AddTask}, so the tabs always see the output of the tasks for that entry.
  \item \code{AtBranch::SetBranchName()} converts the name to a FairRoot branch ID;
        \code{SetBranchId()} notifies only if the ID changed, and remembers the previous ID.
  \item \code{AtSimpleType<T>::Set(value, notify=true)} notifies only if the value changed.
\end{itemize}

On a change of entry, \code{AtViewerManager::Update} calls \code{Exec()} on every tab and then
\code{fPadNum.Notify()}, so pad-level views redraw for the new event even when the pad number did
not change. \code{AtTabMain::Update} shows the pattern used by all tabs: the AtEvent-dependent
parts (pad plane, 3D hits) redraw when \code{fEventBranch} or \code{fEntry} changed, the
pattern-event parts when \code{fPatternEventBranch} or \code{fEntry} changed, and the trace when
\code{fRawEventBranch}, \code{fEntry} or \code{fPadNum} changed.

\paragraph{Tab data: \code{AtTabInfo}.} Each tab owns an \code{AtTabInfo}
(\url{AtEventDisplay/AtTabs/AtTabInfo.h}), a named collection of ``augments'' that give access to
data. \code{AtTabInfoFairRoot<T>} wraps an \code{AtBranch}: its \code{GetInfo()} re-reads the
\code{TClonesArray} from \code{FairRootManager} whenever the selected branch name changed and
returns element 0 cast to \code{T}, or \code{nullptr} when the array is empty. The default name
of an augment is the class name, so a tab retrieves its event with
\code{GetFairRootInfo<AtEvent>()}. \code{AtTabInfoTree} and \code{AtTabInfoBranch<T>}
(\url{AtEventDisplay/AtTabs/AtTabInfoTree.h}) do the same for a non-FairRoot \code{TTree} read
with a \code{TTreeReader}, kept in step with the viewer by observing \code{AtTreeEntry}; the
header recommends creating each tree object once and sharing it through a \code{shared\_ptr}.

\section{Opening a file and browsing events}
\label{sec:vis-open}

\subsection{A minimal launcher}

The following macro is \url{macro/Unpack_HDF5/a1975/run_eve_new.C}, shortened. It opens
\code{run\_NNNN.root} in the current directory, adds a main tab and a two-cell trace tab, and runs
a PSA task on the fly so that the hit threshold can be changed from the sidebar.
\begin{verbatim}
void run_eve_new(int runNum = 11,
                 TString OutputDataFile = "./output.reco_display.root")
{
   TString InputDataFile = TString::Format("./run_%04d.root", runNum);
   TString dir = getenv("VMCWORKDIR");
   TString GeoDataPath = dir + "/geometry/ATTPC_D300torr_v2_geomanager.root";
   TString mapDir = dir + "/scripts/ANL2023.xml";

   FairRunAna *fRun = new FairRunAna();
   fRun->SetSource(new FairFileSource(InputDataFile));
   fRun->SetSink(new FairRootFileSink(OutputDataFile));
   fRun->SetGeomFile(GeoDataPath);
   FairRuntimeDb *rtdb = fRun->GetRuntimeDb();
   rtdb->setFirstInput(new FairParRootFileIo());

   auto fMap = std::make_shared<AtTpcMap>();
   fMap->ParseXMLMap(mapDir.Data());
   auto eveMan = new AtViewerManager(fMap);

   auto tabMain = std::make_unique<AtTabMain>();
   tabMain->SetMultiHit(100);

   auto psa = new AtPSAMax();
   psa->SetThreshold(50);
   AtPSAtask *psaTask = new AtPSAtask(psa);
   psaTask->SetPersistence(false);
   psaTask->SetOutputBranch("AtEventChange");
   auto sidePSA = new AtSidebarPSA(eveMan->GetSidebar());
   sidePSA->SetPSA(psa);
   eveMan->GetSidebar()->AddSidebarFrame(sidePSA);
   eveMan->AddTask(psaTask);

   auto tabPad = std::make_unique<AtTabPad>(2, 1);
   tabPad->DrawRawADC(0, 0);
   tabPad->DrawADC(0, 1);

   eveMan->AddTab(std::move(tabMain));
   eveMan->AddTab(std::move(tabPad));
   eveMan->Init();
}
\end{verbatim}
It is started with \code{root -l 'run\_eve\_new.C(11)'} from a shell where the ATTPCROOT
environment (\code{VMCWORKDIR}, library paths) is set. The pad map file and the geometry file are
experiment specific: map XML files live in \url{scripts/}, geometry ROOT files in
\url{geometry/}. A sink is set because \code{FairRunAna} writes an output file; the macros
name it \code{*.reco\_display.root}.

A launcher that only displays a file that already contains raw and hit branches needs no task;
\url{macro/Unpack_HDF5/a1975/UKF/pipeline/run_eve_UKF.C} is such a case (its header notes that no
parameter file is needed because no PSA runs). Other launchers in the repository include
\url{macro/tests/AT-TPC/run_eve.C}, \url{macro/Simulation/ATTPC/PUMA/run_eve.C} and
\url{macro/Unpack_HDF5/a2091/run_eve.C}.

\subsection{Choosing the branches}

Branches are chosen in the GUI, not in the macro. \code{GenerateBranchLists()} loops over every
branch known to \code{FairRootManager}, keeps those that are a \code{TClonesArray}, steps through
up to the first five entries until the array is non-empty, and files the branch under the class
name of its first element. A branch is therefore offered in the ``Event'' menu whenever it holds
\code{AtEvent} objects, whatever its name (\code{AtEventH}, \code{AtEventCorrected},
\code{AtEventChange}, \ldots). The three drop-downs of \code{AtSidebarBranchControl} list the
branches of type \code{AtRawEvent}, \code{AtEvent} and \code{AtPatternEvent}; each selects
index 0 at start-up, i.e.\ the first branch of that type in FairRoot's branch list. A type with
no branch in the file gets no menu.

\subsection{Navigating}

The basebar holds the event number entry (limited to $[0, N]$ with $N$ the number of chain
entries), Prev and Next buttons, and ``Rerun Event''. Every route ends in
\code{AtTreeEntry::Set}, so every navigation re-runs the attached tasks. ``Rerun Event'' sets the
entry to its current value, which re-runs the tasks on the same event; this is how a parameter
changed in a sidebar add-on (for example the PSA threshold) takes effect. Changing the threshold
alone does not redraw anything. When \code{SetCheckBranch()} has been called, Next and Prev loop
until they reach an entry whose raw event is good.

\section{The tabs}
\label{sec:vis-tabs}

All tabs derive from \code{AtTabBase} (\url{AtEventDisplay/AtTabs/AtTabBase.h}).
\code{AtTabBase::Init()} creates a new window in the right-hand tab widget of the TEve browser
and calls the virtual \code{MakeTab()} (build the GUI) and \code{InitTab()} (register the data
augments). \code{Exec()} is called after the run has processed each entry. Canvas-based tabs
derive from \code{AtTabCanvas}, which provides one \code{TCanvas} divided into
\code{rows}\,$\times$\,\code{cols} pads (the code comment says a \code{TCanvas} is used instead of
\code{TRootEmbeddedCanvas} to keep the toolbar). Table~\ref{tab:vis-tabs} lists the tabs compiled
into the library (\url{AtEventDisplay/CMakeLists.txt}).

\begin{table}[ht]
\centering
\begin{tabular}{p{0.2\textwidth}p{0.17\textwidth}p{0.53\textwidth}}
\toprule
Tab & Base & Shows \\
\midrule
\code{AtTabMain} & \code{AtTabBase} & 3D hits and patterns, pad plane, trace of the selected pad. \\
\code{AtTabPad} & \code{AtTabCanvas} & Grid of trace views of the selected pad. \\
\code{AtTabMacro} & \code{AtTabCanvas} & Grid of pads filled by user lambdas. \\
\code{AtTabBraggCurve} & \code{AtTabMain} & Main view plus energy loss vs range of the first track. \\
\code{AtTabFitted} & \code{AtTabBase} & Overlay of fitted trajectories on the 3D view. \\
\code{AtTabFission} & \code{AtTabMain} & Main view plus uncorrected and corrected fission-fragment hits. \\
\code{AtTabEnergyLoss} & \code{AtTabCanvas} & dE/dx of two fission fragments, sums, fits and ratio. \\
\code{AtTabFF} & \code{AtTabCanvas} & dQ/dZ of two fragments, simulated and experimental. \\
\bottomrule
\end{tabular}
\caption{Tabs available in \code{AtEventDisplay/AtTabs}.}
\label{tab:vis-tabs}
\end{table}

\subsection{\code{AtTabMain}: 3D view, pad plane and trace}

\code{AtTabMain} (\url{AtEventDisplay/AtTabs/AtTabMain.h}) is the tab most launchers start with.
\code{MakeTab} builds a horizontal pack: on the left a TEve 3D viewer with the global and event
scenes; on the right a vertical pack with the pad-trace canvas on top and the pad-plane canvas
below (weight 1.5). When a geometry is loaded (\code{gGeoManager} non-null) the top node is added
to the 3D scene and the volume \code{drift\_volume} is made 80\% transparent.

\paragraph{3D view.} The hits of the selected \code{AtEvent} branch are drawn as a
\code{TEvePointSet} named ``Hits'' inside a \code{TEveEventManager} called ``AtEvent''. The hits
of each track in the selected \code{AtPatternEvent} are drawn as separate point sets
\code{Track\_0}, \code{Track\_1}, \ldots{} with a fixed colour list, and the noise hits of the
pattern event as a set named ``Noise''. With \code{SetDrawProjection(true)} the projection of
each track's \code{AtPattern} is also drawn on the pad plane.

\paragraph{Pad plane.} The pad plane is the \code{TH2Poly} returned by
\code{AtMap::GetPadPlane()}, drawn with option \code{"COL L0"} and filled with the charge of
each hit at its $(x,y)$ position. A hit is skipped when its charge is below the threshold or when
its pad has more hits than the multi-hit limit.

\paragraph{Pad selection and trace.} A left click on the pad plane runs the
static method \code{AtTabMain::SelectPad()} (registered with \code{TCanvas::AddExec}). It converts the click to
a \code{TH2Poly} bin, maps the bin to a pad number with \code{AtMap::BinToPad}, and sets
\code{fPadNum}. Every observer of \code{fPadNum} reacts: \code{AtTabMain} draws the ADC of that pad
from the selected \code{AtRawEvent} branch (bins 0--511), \code{AtTabPad} redraws its grid, and
\code{AtSidebarPadControl} shows the pad number and its electronics reference. A pad can also be
typed into the sidebar entry and redrawn with ``Redraw Pad''.

\paragraph{Setters} (defaults from the header):
\begin{center}
\begin{tabular}{p{0.38\textwidth}p{0.14\textwidth}p{0.38\textwidth}}
\toprule
Setter & Default & Meaning \\
\midrule
\code{SetThreshold(Int\_t)} & 0 & Minimum hit charge drawn. \\
\code{SetMultiHit(Int\_t)} & 10 & Maximum hits in a pad for its hits to be drawn. \\
\code{SetDrawProjection(Bool\_t)} & \code{false} & Draw pattern projections on the pad plane. \\
\code{SetHitAttributes(TAttMarker)} & pink, full dot medium, size 1 & Hit marker. \\
\code{DumpEvent(std::string)} & --- & Write the hits of the current event to a text file. \\
\bottomrule
\end{tabular}
\end{center}
The default multi-hit limit of 10 hides every pad with more than 10 hits; launchers that use PSA
methods producing many hits per pad raise it (all a1975 launchers use \code{SetMultiHit(100)}).
Note that \url{docs/subsystems/visualization.md} quotes 100 in its example, which is the value
set in the macros, not the class default.

\subsection{\code{AtTabPad}: trace grid}

\code{AtTabPad(nRow = 1, nCol = 1, name = "AtPad")} (\url{AtEventDisplay/AtTabs/AtTabPad.h})
divides its canvas into a grid; each cell is assigned one kind of trace of the selected pad. The
cell is given as \code{(row, col)}:
\begin{center}
\begin{tabular}{p{0.38\textwidth}p{0.52\textwidth}}
\toprule
Call & Cell content \\
\midrule
\code{DrawADC(row, col)} & Pedestal-subtracted ADC of the pad. \\
\code{DrawRawADC(row, col)} & Raw ADC. \\
\code{DrawArrayAug(augName, row, col)} & A named array augment of the pad (e.g.\ a deconvolved
  charge). \\
\code{DrawAuxADC(auxName, row, col)} & An auxiliary pad by name (e.g.\ \code{"IC"}); independent of
  the selected pad. \\
\code{DrawFPN(row, col)} & Fixed-pattern-noise trace. \\
\code{DrawGenTrace(id, row, col)}, \code{DrawRawGenTrace(id, row, col)} & A generic trace by ID. \\
\code{DrawHits(row, col)} & Overlay the hits of that pad on the cell: a Gaussian of integral $Q$
  when the hit has a $z$ variance, a point otherwise. \\
\bottomrule
\end{tabular}
\end{center}
The header states that \code{DrawHits} requires an \code{AtDigiPar} parameter container in the
runtime database. The tab reads the selected \code{AtEvent} and \code{AtRawEvent} branches.

\subsection{\code{AtTabMacro}: user-defined plots}

\code{AtTabMacro(nRow = 1, nCol = 1, name = "Macro")} fills its pads by calling user functions
at three moments:
\begin{itemize}
  \item \code{SetDrawTreeFunction(f, row, col)}: called once, when the tab is built;
  \item \code{SetDrawEventFunction(f, row, col)}: called from \code{Exec()}, after each entry;
  \item \code{SetDrawPadFunction(f, row, col)}: called when the selected pad changes, with the
        pad number as second argument.
\end{itemize}
The functions receive the tab's \code{AtTabInfo*}. The tab registers no augment itself; the macro
adds what the functions need. \url{macro/e12014/adam/unpack/run_eve_species.C} adds a HiRAEVT
tree for a PID plot:
\begin{verbatim}
auto heistTree = std::make_shared<AtTabInfoTree>(
   "E12014", evtInputDataFile, AtViewerManager::Instance()->GetCurrentEntry());
auto heistInfo =
   std::make_shared<AtTabInfoBranch<HTMusicIC>>(heistTree, "MUSIC");
auto tabMac = std::make_unique<AtTabMacro>();
tabMac->GetTabInfo()->AddAugment(std::move(heistInfo), "MusicIC");
tabMac->GetTabInfo()->AddAugment(std::move(heistTree));
tabMac->SetDrawTreeFunction(PlotPID);
\end{verbatim}
The class comment describes it as ``probably needs refactor''.

\subsection{\code{AtTabFitted}: fitted trajectories}

\code{AtTabFitted(name = "Fitted", branchName = "AtTrackingEvent")}
(\url{AtEventDisplay/AtTabs/AtTabFitted.h}) builds no window of its own; it adds an
``AtTrackingEvent'' element to the 3D scene of \code{AtTabMain} and draws one \code{TEveLine}
per \code{AtFittedTrack}, plus a point set for the fitted vertex. It looks the branch up by name
in \code{FairRootManager} instead of going through an \code{AtBranch}; the header explains that
the \code{AtBranch} route resolved to the wrong branch with chained sources. Two drawing modes
exist:
\begin{itemize}
  \item default: an analytical helix sampled from the fit kinematics
        (\code{GetKinematics()}) in the field \code{SetBField(Bz)} (default 4.0~T, the PUMA value),
        up to \code{SetHelixSampling(maxArc\_mm, nSamples)} (defaults 500~mm and 200 points);
  \item \code{SetDrawSmoothed(true)}: the vertex followed by the stored smoothed positions. The
        header calls this ``the honest per-filter trajectory'', valid for both UKF and genfit.
\end{itemize}
Frame options: \code{SetZPadPlane(z)} (default 0) maps the lab-frame $z$ stored by the fitters
back to the digitisation frame as $z_\mathrm{digi} = z_\mathrm{pad} - z_\mathrm{lab}$ in
smoothed mode, so the line overlays the hits; with 0, \code{SetFlipZ(bool)} multiplies $z$ by
$\pm1$. Styling: \code{SetTrackColor} (negative keeps a per-track palette),
\code{SetLineStyle} (default 1), \code{SetLineWidth} (default 2), \code{SetDrawVertex} (default
true). Two instances pointed at two branches overlay two fitters, as in
\url{macro/Unpack_HDF5/a1975/UKF/pipeline/run_eve_both_a1975.C}:
\begin{verbatim}
FairFileSource *source = new FairFileSource(recoFile);
source->AddFriend(fitFile); // merge the two fit branches onto the same entries
...
auto tabGenfit =
   std::make_unique<AtTabFitted>("Genfit", "AtTrackingEventGenfit");
tabGenfit->SetDrawSmoothed(true);
tabGenfit->SetZPadPlane(1000.0);
tabGenfit->SetTrackColor(kRed + 1);
auto tabUKF = std::make_unique<AtTabFitted>("UKF", "AtTrackingEventUKF");
tabUKF->SetDrawSmoothed(true);
tabUKF->SetZPadPlane(1000.0);
tabUKF->SetTrackColor(kAzure + 1);
tabUKF->SetLineStyle(2);
\end{verbatim}
The macro header warns that the reco file and the friend file must be entry-aligned, i.e.\ the fit
file must have been produced from the same run over the events being browsed.

\subsection{Other tabs}

\begin{itemize}
  \item \code{AtTabBraggCurve} extends \code{AtTabMain} with a canvas of energy loss versus range,
        filled from \code{AtTrack::GetBraggCurve()} of the \emph{first} track (index 0) of the
        selected pattern event; with no track the histogram is empty.
  \item \code{AtTabFission} extends \code{AtTabMain} with an ``AtFissionEvent'' element holding
        uncorrected hits and three sets of corrected hits. It owns its own \code{AtBranch}, returned
        by \code{GetFissionBranch()}, which is passed to the two fission analysis tabs below.
  \item \code{AtTabEnergyLoss(AtBranch\&)} computes dE/dx curves of the two fragments of an
        \code{AtFissionEvent}: stacked curves along the track and along $z$, charge-sum and
        Gaussian-fit dQ/dZ, their ratio, and a proxy for the fragment charge.
  \item \code{AtTabFF(AtBranch\&, plotADC = false)} compares simulated and experimental dQ/dZ of the
        two fragments.
\end{itemize}
These four tabs were written for specific experiments (e12014 fission); their exact physics
content is defined only by the source files in \url{AtEventDisplay/AtTabs/}.

\section{The sidebar and its add-ons}
\label{sec:vis-sidebar}

\code{AtEventSidebar} (\url{AtEventDisplay/AtSidebar/AtEventSidebar.h}) is a \code{TGMainFrame}
holding a list of \code{AtSidebarFrame}s. Frames are added with \code{AddSidebarFrame()} before
\code{Init()}; their widgets are only created in \code{FillFrame()}, which \code{Init()} calls
after the run is initialised. The fixed frames were listed in Section~\ref{sec:vis-viewer}.
\code{UsePictureButtons(bool)} switches Prev/Next between picture and text buttons; the source
comment says picture buttons ``break some machines'', and the default is text buttons.

User add-ons derive from \code{AtSidebarAddon}, which provides \code{AddInfoBox(label)} (a text
field updated with \code{SetInfoString}) and \code{AddIntBox(label, function, min, max)} (an
integer entry that calls a slot of the derived class). Available add-ons:
\begin{center}
\begin{tabular}{p{0.3\textwidth}p{0.6\textwidth}}
\toprule
Add-on & Controls \\
\midrule
\code{AtSidebarPSA} & ``Threshold'' (0--4000) of the \code{AtPSA} given with \code{SetPSA()};
  initialised from \code{AtPSA::GetThreshold()}. \\
\code{AtSidebarPSADeconv} & Threshold plus ``Filter Order'' and ``Filter Cutoff'' of a
  deconvolution PSA. \\
\code{AtSidebarPSAIterDeconv}, \code{AtSidebarPSATBAvg} & Parameters of the iterative-deconvolution
  and time-bucket-averaging PSA (see their headers in \url{AtEventDisplay/AtSidebar/PSA/}). \\
\code{AtSidebarInfoMacro} & Displays the string returned by a user \code{std::function<std::string()>};
  refreshed on every entry change. \\
\code{AtSidebarIntMacro} & An integer entry, bounded by \code{SetBounds(min, max)}, whose value is
  passed to a user \code{std::function<void(int)>}. \\
\bottomrule
\end{tabular}
\end{center}
The PSA add-ons change the PSA object in place; the change reaches the display after ``Rerun
Event'' re-executes the \code{AtPSAtask} (Section~\ref{sec:vis-open}). Pointing the PSA task at a
new output branch (\code{"AtEventChange"} in the example) and selecting that branch in the
``Event'' menu shows the re-analysed hits next to the stored ones.

\section{The \code{Deprecated/} directory}
\label{sec:vis-deprecated}

The directory \url{AtEventDisplay/Deprecated/} holds the previous
viewer: the event managers \code{AtEventManager}, \code{AtEventManagerProto} and
\code{AtEventManagerS800}, and the draw tasks \code{AtEventDrawTask},
\code{AtEventDrawTaskProto} and \code{AtEventDrawTaskS800}. In that design the event manager is a
\code{TEveEventManager} and drawing is done by a FairTask whose branch names are fixed in the
macro (\code{SetRawEventBranch}, \code{SetEventBranch}, \code{SetMultiHit},
\code{Set3DHitStyleBox}, \ldots). These classes are still listed in \code{SRCS} of
\url{AtEventDisplay/CMakeLists.txt} and in the LinkDef, so they are compiled and loadable, and
older macros still use them (for example \url{macro/e12014/run_eve.C} and
\url{macro/e15250/run_eve.C}). Nothing in the code marks them as deprecated other than the
directory name, and no removal is documented. New launchers should use \code{AtViewerManager};
the subsystem notes in \url{docs/subsystems/visualization.md} describe only the new API.

Some macros mix the two generations or call methods that no longer exist: for instance
\url{macro/e12014/curtis/viewer/run_eve_dQdtA.C} calls \code{eveMan->GetRawEventName()}, which
\code{AtViewerManager} does not provide. Macros under \url{macro/Simulation/ATTPC/16C_pp/} also
contain \code{eventDisplay.C}, which uses FairRoot's own \code{FairEventManager}.

\section{Fitter GUIs}
\label{sec:vis-fitter}

\subsection{Design note and \code{AtUKFDisplay}}

\url{docs/development/FitterGUI.md} is a design document for an interactive UKF fitter display
modelled on GenFit's \code{EventDisplay}: a singleton with event navigation, live refitting from
GUI parameters, and display flags (\code{C} clusters, \code{H} raw hits, \code{F} filtered,
\code{S} smoothed, \code{R} residuals, \code{E} covariance ellipses, \code{M} MC truth). It plans
three phases: a plotting macro, a compiled class, and advanced features (MC overlay, error bands,
UKF/GenFit comparison, batch mode). Phase 1 exists as
\url{macro/Simulation/ATTPC/16C_pp/display_ukf.C}. Phase 2 exists as \code{AtUKFDisplay}, but in
\url{AtReconstruction/AtFitter/AtUKFDisplay.h} (library \code{AtReconstruction}), not in the
\code{AtEventDisplay/AtUKFDisplay/} directory the document proposes, and it differs from the
design in that it uses only \code{TCanvas} (\code{fUseEve} is forced to false; no OpenGL).

The launcher is \url{macro/Simulation/ATTPC/16C_pp/run_ukf_display.C}:
\begin{verbatim}
void run_ukf_display(int startEvent = 1)
{
   FairLogger::GetLogger()->SetLogScreenLevel("WARNING");
   auto *display = AtUKFDisplay::GetInstance();
   display->LoadFiles("data/output_digi.root", "data/output_ukf_only.root",
                      "data/attpcsim.root");
   display->GotoEvent(startEvent);
}
\end{verbatim}
\code{LoadFiles(digi, fitted = nullptr, mc = nullptr)} reads the tree \code{cbmsim} and the
branches \code{AtPatternEvent} (required), \code{AtTrackingEvent} and \code{MCTrack} (optional).
The GUI consists of a control frame and two canvases: \code{cUKFTrack} (XY, XZ, YZ projections of
clusters, optional raw hits, and the fit) and \code{cUKFDiag} (cluster charge, residuals
$|$smoothed$-$measured$|$, $z_\mathrm{lab}$ per cluster, and a text box with $p$, KE, $\theta$,
$\phi$). The controls and their initial values, read in \code{CreateFitter()} each time a fit is
requested:
\begin{center}
\begin{tabular}{p{0.34\textwidth}p{0.18\textwidth}p{0.38\textwidth}}
\toprule
Control & Initial value & Passed to \code{AtFitterUKF} \\
\midrule
Particle & Proton & mass and charge (proton, deuteron, alpha, $\pi^\pm$) \\
UKF alpha & $10^{-3}$ & \code{SetUKFParameters(alpha, 2.0, 0.0)} \\
Measurement sigma & 2.0 & \code{SetMeasurementSigma} \\
Min clusters & 10 & \code{SetMinClusters} \\
E-loss scale & 1.0 & read but not passed to the fitter in \code{CreateFitter()} \\
ZPadPlane & 1000.0 & \code{SetZPadPlane} \\
Momentum seed & 0 (use PR \brho) & \code{SetMomentumSeed} if $>0$ \\
Momentum sigma fraction & 0.3 & \code{SetMomentumSigmaFrac} \\
$B_z$ [T] & 2.85 & \code{SetBField} \\
Gas density [g/cm$^3$] & $3.553\times10^{-5}$ & \code{AtELossCATIMA} \\
Energy straggling, per-cluster covariance & off & corresponding setters \\
\bottomrule
\end{tabular}
\end{center}
``Re-cluster + Fit'' re-clusters the displayed track with Smooth3D (radius 10, distance 20) or
GroupByN (15 hits per cluster) before fitting. ``Use MC truth seed'' replaces the seed momentum by
that of the first MC track with PDG code 2212. ``Auto-fit on navigate'' fits on each event change
when no stored fit exists. \code{SetParticle(kind)}, \code{SetBField(T)} and
\code{SetGasDensity(g/cm$^3$)} preset the GUI from a macro after \code{LoadFiles()}.

\subsection{\code{gui\_fit\_dp.C}: live genfit refit of one track}

\url{macro/Unpack_HDF5/a1975/D2_UKF/gui_fit_dp.C} is an analysis-specific GUI (a1975,
16C(d,p)) that calls \code{AtGenfitter} on one stored track with parameters from the window, about
one second per fit according to its header. It reads
\code{<gfDir>/<runTag>\_genfitter\_p.root}, picks the \code{AtTrackingEvent} of the given entry
and the track with the given ID, and also shows the fit stored in the production for comparison.
\begin{verbatim}
root -l 'gui_fit_dp.C("run_0020_multifit",9789,0)'
\end{verbatim}
The parameters are the cluster ordering (production $z$-sort, stored cluster order, arc walk,
vertex-first), measurement sigma (initial 4.0), minimum and maximum iterations (2 and 5), a cluster
index window, and switches for material effects, CATIMA, the backward-seed fix and a Spyral seed.
Clusters can be selected by dragging a box in a projection (keep or drop), or by a charge cut;
``find'' searches the longest prefix of the track whose $\chi^2/\mathrm{ndf}$ is below a limit
(initial 0.1). The canvas has three projections and a fourth pane with the RANSAC line from
pattern recognition; ``PNG'' saves it under \code{plots/}. The header records a regression check:
at default settings, run\_0020 entry 9789 must give KE 0.385~MeV, $\theta$ 96.92$^\circ$,
$\chi^2/\mathrm{ndf}$ 143.66, and the banner warns if it does not. The default paths point to
\url{/mnt/f/a1975/}, so the macro needs arguments on any other machine.

\section{Other interactive and visual macros}
\label{sec:vis-macros}

\subsection{PID gate drawers}

Polygon gates are stored as \code{AtTools::AtCut2D} (\url{AtTools/AtPID/AtCut2D.h}): a name,
the two axis names, and the vertex list, read and written as JSON in the spyral\_utils
\code{Cut2D} format (\code{LoadJSON}, \code{WriteJSON}, \code{IsInside}).
\code{AtTools::AtParticleID} couples a cut with $Z$ and $A$. The gates are drawn interactively by
experiment-specific macros, all of the form \code{gate\_draw\_*.C} or \code{draw\_gate\_*.C}:
\begin{itemize}
  \item a1975: \url{macro/Unpack_HDF5/a1975/D2_UKF/pid/gate_draw_dp.C},
        \code{gate\_draw\_dt.C}, \code{draw\_gate\_dt.C}, \code{pid\_draw\_gate\_d2.C} (same
        directory), and \url{macro/Unpack_HDF5/a1975/UKF/pid/gate_draw_pp.C},
        \code{pid\_draw\_gate\_a1975.C};
  \item other experiments: \url{macro/Unpack_HDF5/C15p/pid/gate_draw_C15p.C},
        \url{macro/Unpack_HDF5/C15d/pid/gate_draw_C15d.C},
        \url{macro/Unpack_HDF5/a1954_Be12/UKF/pid/gate_draw_pt_Be12.C},
        \url{macro/Unpack_HDF5/a2091/UKF/pid/draw_gate_C15.C};
  \item gate editors: \code{pid\_gate\_edit.C} and \code{ex\_gate\_edit.C} in
        \url{macro/Unpack_HDF5/a1975/UKF/pp/}.
\end{itemize}

\paragraph{Example: \code{gate\_draw\_dp.C}.} The drawer reads a cached PID plane (a tree
\code{pts} with \sqdedx, \brho, polar angle and IC amplitude, built by \code{pid\_plane\_dt.C})
and opens a \code{TGMainFrame} with an embedded canvas and the buttons ``Draw new gate'',
``Evaluate'', ``Locus check'', ``Locus dots on/off'', ``(d,d) band on/off'', ``Save JSON'',
``Save PNG'', ``Redraw'' and ``Quit'', plus entries for axis ranges, binning and the IC window.
``Draw new gate'' waits for a \code{TCutG} (\code{WaitPrimitive("CUTG")}): click the vertices
and double-click to close. ``Locus check'' plots \brho{} against polar angle for the tracks inside
the polygon with the two-body kinematic curve for the chosen beam energy overlaid. The defaults of
the function signature are output \code{pid/proton\_d2\_dv1104.json}, cache
\code{pid/pid\_plane\_dt\_dv1104.root}, $x_\mathrm{max}=55$, $y_\mathrm{max}=1.6$, IC window
900--1300, $E_\mathrm{beam}=184.25$~MeV, $Z=1$, $A=1$. Save behaviour that matters:
\begin{itemize}
  \item the file name follows the name typed in the GUI, and a species keyword in that name
        (\code{proton}, \code{deuteron}, \code{triton}, \code{3he}, \code{alpha}, \code{4he})
        overrides $Z$ and $A$;
  \item an existing file is copied to \code{.bak} before being replaced, with a warning.
\end{itemize}
Launch from the analysis directory, either at an interactive prompt or detached with stdin kept
open (from \url{macro/Unpack_HDF5/a1975/D2_UKF/pid/run_gate_dp.C}):
\begin{verbatim}
cd macro/Unpack_HDF5/a1975/D2_UKF
root -l 'pid/run_gate_dp.C'
setsid bash -c 'tail -f /dev/null | root -l pid/run_gate_dp.C' &
\end{verbatim}

\subsection{Event steppers}

\url{macro/Unpack_HDF5/a1975/UKF/pipeline/cycleEventsGui.C} opens a genfitter output file
(branches \code{AtTrackingEvent} and \code{AtPIDEvent}), loads an \code{AtParticleID} gate, and
shows XY, ZX and ZY projections with Prev, Next, Goto, Save PNG and a ``gated-only'' check box.
Tracks inside the gate are drawn in colour, the others in grey. It uses a 2D canvas only and works
over X11 without OpenGL. \code{cycleEvents.C} in the same directory does the same with terminal
commands (Enter, \code{p}, a number, \code{s}, \code{g}, \code{q}). Both print a warning when the
gate file does not load, because ``gated-only'' then shows no events.

\subsection{Batch (headless) displays}
\label{sec:vis-batch}

These macros write PNG files and are meant for \code{root -b -q}:
\begin{itemize}
  \item \url{macro/Unpack_HDF5/a1975/UKF/pipeline/event_png_UKF.C}: one event of the
        \code{AtEventH} branch as XY, XZ, YZ and a software-rendered 3D scatter
        (\code{TGraph2D}); event $-1$ picks the event with most hits among the first
        \code{scanN} (default 200).
  \item \url{macro/Unpack_HDF5/a1975/UKF/pipeline/event_png_batch_UKF.C}: a contact sheet of
        XY views over an entry range, plus one detailed PNG per event, written to
        \code{<file>\_pngs/}.
  \item \url{macro/Unpack_HDF5/a1975/UKF/pipeline/contactSheet.C}: a grid of XY thumbnails of
        fitted (optionally gated) events labelled with entry, $\theta$, KE and
        $\chi^2/\mathrm{ndf}$.
\end{itemize}
\begin{verbatim}
root -b -q 'event_png_UKF.C("run_0116_disp", -1)'
root -b -q 'pipeline/contactSheet.C("/tmp/run_0106_genfitter_gated.root",
            "pid/deuteron_band.json","/tmp/contact.png",true,6)'
\end{verbatim}
(The second command is one line in practice.) Any macro that draws on a \code{TCanvas} can be
used the same way by adding \code{c->SaveAs("file.png")} and running with \code{-b}.

\section{Pitfalls}
\label{sec:vis-pitfalls}

\begin{itemize}
  \item \textbf{Order of construction.} \code{FairRunAna} first, then \code{AtViewerManager},
        then tabs, then \code{Init()}. Tabs created before the viewer throw; a second viewer
        aborts the process.
  \item \textbf{Hidden hits.} The \code{AtTabMain} multi-hit limit defaults to 10; pads with more
        hits are not drawn in 3D or on the pad plane.
  \item \textbf{Geometry volume name.} \code{AtTabMain::MakeTab} looks up the volume
        \code{drift\_volume} with \code{gGeoManager->GetVolume} and sets its transparency
        without checking the result; a geometry file without a volume of that name dereferences a null pointer.
  \item \textbf{Trace length.} \code{AtTabMain::DrawWave} copies 512 ADC samples.
  \item \textbf{Branch type detection.} A branch that is empty in the first five entries is not
        offered in any menu (``Failed to find type of branch'' in the log).
  \item \textbf{Parameter changes need a rerun.} Sidebar PSA controls change the PSA object only;
        press ``Rerun Event''.
  \item \textbf{Friend files.} With \code{AddFriend}, the friend must be entry-aligned with the
        main file, or the fit drawn belongs to another event.
  \item \textbf{\code{AtUKFDisplay}.} \code{SetOptions()} stores a flag string that the drawing
        code does not read. The display shows the track with most hits but takes the stored fit
        from the first fitted track of the event, which need not be the same track. The CATIMA
        model in \code{CreateFitter()} is built with projectile $Z=1$ and a pure hydrogen
        material for every particle choice, including the alpha.
  \item \textbf{Exiting GUIs.} The a1975 GUIs (\code{gui\_fit\_dp.C}, the gate drawers) exit with
        \code{gSystem->Exit(0)} and their launchers do not call \code{gApplication->Run()}; the
        header of \code{run\_gate\_dp.C} explains that the earlier combination of a nested
        \code{Run()} and \code{Terminate()} re-executed the macro and reopened the window about
        twice per second.
  \item \textbf{Stale comments.} Some drawer comments were copied from the macro they were ported
        from: the \code{@param} block of \code{gate\_draw\_dp.C} still describes a 14C beam and
        \code{pid\_plane\_pt\_C14.C}. Trust the function signature for defaults.
\end{itemize}
